% !TeX TS-program = pdflatex
\documentclass[11pt,a4paper]{article}
\usepackage[margin=0.9in,headheight=14pt]{geometry}
\usepackage[T1]{fontenc}
\usepackage{lmodern}
\usepackage{amsmath,amssymb}
\usepackage{enumitem,booktabs,fancyhdr}
\setlength{\parindent}{0pt}
\setlength{\parskip}{0.5em}
\linespread{1.08}
\pagestyle{fancy}
\fancyhf{}
\fancyhead[C]{\small Probability \& Statistics for EECS: Homework \#04}
\fancyfoot[C]{\thepage}
\renewcommand{\headrulewidth}{0.4pt}
\renewcommand{\footrulewidth}{0.4pt}
\setlist[enumerate]{label=(\alph*),leftmargin=*,itemsep=1em}
\newcommand{\Prob}{P}
\title{SI 140A.02 Probability \& Statistics for EECS, Fall 2026\\[0.4em]Homework 4}
\author{%
\begin{tabular}{c}
Name: \underline{\makebox[2.6in][c]{Yiheng Zhang}}\\[0.6em]
Student ID: \underline{\makebox[2.6in][c]{2025531068}}
\end{tabular}}
\date{
Due date: Due on Oct. 18, 2026, 23:59 UTC+8}
\begin{document}
\maketitle
\thispagestyle{plain}
\vspace{1em}
There are five problems in this assignment.

Read all the instructions below carefully before you start working on the assignment, and before you make a submission.
\begin{itemize}
\item You are required to write down all the major steps towards making your conclusions; otherwise you may receive limited credit for the problem.
\item Write your homework in English; otherwise you will receive no credit for this homework.
\item Any form of plagiarism will result in no credit for this homework.
\item You are required to submit your homework in PDF format generated by LaTeX, otherwise you will receive no credit for this homework.
\end{itemize}

\clearpage
\section*{Problem 1}
A token starts at position $1$ on the four positions $0,1,2,3$. It stops when it first reaches $0$ or $3$. Fix $r\in[0,1]$. All random choices in different rounds are independent. In an interior position, either of the following rules can be used throughout the process:
\begin{itemize}
\item \textbf{Pause rule:} With probability $r$, the token stays where it is; with probability $1-r$, a fresh fair coin moves it one position left or right.
\item \textbf{Restart rule:} With probability $r$, the token is placed at position $1$; with probability $1-r$, a fresh fair coin moves it one position left or right.
\end{itemize}
\begin{enumerate}
\item Under the pause rule, determine the probability of ever reaching position $3$ for every $r\in[0,1]$. Justify your answer.
\item Answer part (a) under the restart rule.
\item A student claims that pausing and restarting have the same effect because both only delay the token. Determine exactly when the two success probabilities are equal. For $r<1$, let $T$ be the number of rounds until the token stops. Find $\Prob(T>n)$ for every integer $n\geq0$ under each rule, and use your result to prove that each process stops with probability $1$. \emph{Hint:} Condition on either possible interior position after $n$ rounds.
\end{enumerate}

\subsection*{Solution}

Let $q_k$ be the probability of ever reaching position $3$ when the token
starts at position $k$. The endpoints are absorbing, so $q_0=0$ and $q_3=1$.
We condition on the outcome of the first round.

\begin{enumerate}
\item \textbf{Pause rule.} With probability $r$ the token stays where it is,
and with probability $(1-r)/2$ it moves to each neighbour. Hence
\begin{align*}
q_1 &= r\,q_1 + \frac{1-r}{2}\,q_0 + \frac{1-r}{2}\,q_2,\\
q_2 &= r\,q_2 + \frac{1-r}{2}\,q_1 + \frac{1-r}{2}\,q_3 .
\end{align*}
Using $q_0=0$ and $q_3=1$,
\begin{align*}
q_1 &= r\,q_1 + \frac{1-r}{2}\,q_2,\\
q_2 &= r\,q_2 + \frac{1-r}{2}\,q_1 + \frac{1-r}{2}.
\end{align*}

If $r=1$ the token never moves, so starting from position $1$ it never
reaches $3$; hence $q_1=0$.

Assume now $0\le r<1$. The first equation gives
\[
(1-r)q_1=\frac{1-r}{2}\,q_2 \quad\Longrightarrow\quad q_1=\frac{q_2}{2},
\]
and the second becomes
\[
(1-r)q_2=\frac{1-r}{2}\,q_1+\frac{1-r}{2}
\quad\Longrightarrow\quad q_2=\frac{q_1}{2}+\frac12 .
\]
Substituting $q_1=q_2/2$: $\;q_2=\tfrac{q_2}{4}+\tfrac12$, so
$\tfrac34 q_2=\tfrac12$, hence $q_2=\tfrac23$ and $q_1=\tfrac13$.

Therefore, under the pause rule,
\[
\Prob(\text{reaching position }3)=
\begin{cases}
\dfrac13, & 0\le r<1,\\[2mm]
0, & r=1.
\end{cases}
\]

\item \textbf{Restart rule.} With probability $r$ the token is placed back at
position $1$, and with probability $(1-r)/2$ it moves to each neighbour:
\begin{align*}
q_1 &= r\,q_1 + \frac{1-r}{2}\,q_0 + \frac{1-r}{2}\,q_2,\\
q_2 &= r\,q_1 + \frac{1-r}{2}\,q_1 + \frac{1-r}{2}\,q_3 .
\end{align*}
With $q_0=0$ and $q_3=1$,
\begin{align*}
q_1 &= r\,q_1 + \frac{1-r}{2}\,q_2,\\
q_2 &= \Bigl(r+\frac{1-r}{2}\Bigr)q_1 + \frac{1-r}{2}
     = \frac{1+r}{2}\,q_1 + \frac{1-r}{2}.
\end{align*}

If $r=1$ the token is reset to $1$ every round, so it never reaches $3$;
again $q_1=0$.

Assume $0\le r<1$. As before the first equation yields $q_1=q_2/2$, and
substituting into the second,
\[
q_2=\frac{1+r}{2}\cdot\frac{q_2}{2}+\frac{1-r}{2}
\quad\Longrightarrow\quad
q_2\Bigl(1-\frac{1+r}{4}\Bigr)=\frac{1-r}{2}
\quad\Longrightarrow\quad
q_2\,\frac{3-r}{4}=\frac{1-r}{2},
\]
so
\[
q_2=\frac{2(1-r)}{3-r},\qquad q_1=\frac{1-r}{3-r},
\]
which also equals $0$ at $r=1$. Thus under the restart rule
\[
\Prob(\text{reaching position }3)=\frac{1-r}{3-r}\qquad (0\le r\le 1).
\]

\item Parts (a) and (b) give
\[
\Prob(\text{reach }3)=
\begin{cases}
\dfrac13, & 0\le r<1,\\[2mm]
0, & r=1,
\end{cases}
\quad\text{(pause)}
\qquad
\Prob(\text{reach }3)=\frac{1-r}{3-r}
\quad\text{(restart)} .
\]
For $0<r<1$ the equation $\frac{1-r}{3-r}=\frac13$ would force
$3(1-r)=3-r$, i.e.\ $r=0$, which is impossible; so on $0<r<1$ the two
probabilities differ. At the endpoints they agree: for $r=0$ both equal
$1/3$, and for $r=1$ both equal $0$. Hence the two success probabilities
coincide \emph{exactly} when
\[
r=0 \quad\text{or}\quad r=1 .
\]
Thus for $0<r<1$ pausing and restarting do \emph{not} have the same effect.
Restarting from position $2$ moves the token back to $1$, whereas pausing
leaves it at $2$, so restarting can change which endpoint is reached.

Next, let $X_n$ be the token's position
after $n$ rounds; the token has not stopped by then precisely when it still
occupies an interior position, so
\[
\{T>n\}=\{X_n\in\{1,2\}\}.
\]
Write $a_n=\Prob(X_n=1)$ and $b_n=\Prob(X_n=2)$; since the token starts at $1$,
\[
a_0=1,\qquad b_0=0 .
\]

Under the \emph{pause rule}, conditioning on the $(n-1)$-st position,
\begin{align*}
a_n &= r\,a_{n-1}+\frac{1-r}{2}\,b_{n-1},\\
b_n &= \frac{1-r}{2}\,a_{n-1}+r\,b_{n-1},
\end{align*}
so
\[
a_n+b_n=\frac{1+r}{2}\,(a_{n-1}+b_{n-1}).
\]

Under the \emph{restart rule},
\begin{align*}
a_n &= r\,a_{n-1}+\frac{1+r}{2}\,b_{n-1},\\
b_n &= \frac{1-r}{2}\,a_{n-1},
\end{align*}
and again
\[
a_n+b_n=\frac{1+r}{2}\,(a_{n-1}+b_{n-1}).
\]

The same recurrence for $a_n+b_n$ therefore holds under either rule.
Iterating it,
\[
\Prob(T>n)=a_n+b_n=\Bigl(\frac{1+r}{2}\Bigr)^{n}(a_0+b_0)
=\Bigl(\frac{1+r}{2}\Bigr)^{n},
\qquad n\ge0 .
\]
The two rules have the same stopping-time distribution, but this does not
imply the same probability of reaching each endpoint: restarting changes
the token's position and hence can change its route to absorption.

Finally, the events
$\{T>n\}$ are decreasing, so by continuity of probability from above,
\[
\Prob(T=\infty)=\Prob\Bigl(\bigcap_{n\ge0}\{T>n\}\Bigr)
=\lim_{n\to\infty}\Prob(T>n).
\]
For $r<1$ we have $\frac{1+r}{2}<1$, hence
$\Prob(T>n)=\bigl(\frac{1+r}{2}\bigr)^{n}\to0$ and therefore
\[
\Prob(T=\infty)=0 .
\]
So under each rule the token stops with probability $1$.
\end{enumerate}

\clearpage
\section*{Problem 2}
Two algorithms, $A$ and $B$, are each tested on $100$ tasks. Across the $200$ test records, exactly $100$ tasks are easy and $100$ are hard. Their exact success rates are:
\begin{center}
\begin{tabular}{ccc}
\toprule Algorithm & Easy tasks & Hard tasks\\\midrule
$A$ & $4/5$ & $2/5$\\
$B$ & $9/10$ & $3/5$\\\bottomrule
\end{tabular}
\end{center}
Each algorithm's easy/hard versus success/failure table has strictly positive integer entries. Let $x$ be the number of easy tasks assigned to $A$.
\begin{enumerate}
\item Find every feasible $x$ for which $A$ has a strictly greater overall success rate than $B$. For the smallest such $x$, give both complete $2\times2$ count tables.
\item From this part onward, use a probability model without finite-sample or integer-count constraints. The displayed within-type success probabilities remain fixed. Both algorithms receive the same easy-task proportion $t\in[0,1]$. An audit reports that their overall success probabilities differ by $13/100$, but omits which algorithm performs better. Recover $t$, identify the better algorithm, and calculate both overall success probabilities.
\item Now let $s,t\in(0,1)$ be the easy-task proportions for $A,B$, respectively, so each algorithm receives both task types. Construct one pair $(s,t)$ satisfying $|s-t|\leq2/5$ for which $A$'s overall success probability is exactly $1/100$ greater than $B$'s. Calculate both probabilities and check all the constraints for your pair. Briefly explain how this is compatible with $B$'s superiority within each task type.
\end{enumerate}

\subsection*{Solution}

Let $x$ be the number of easy tasks assigned to $A$. Since each algorithm
receives $100$ tasks, $A$ gets $x$ easy and $100-x$ hard tasks, while $B$
gets $100-x$ easy and $x$ hard tasks.

\begin{enumerate}
\item By the law of total probability (conditioning on the task type), the
overall success probability of $A$ is
\[
\Prob_A=\frac45\cdot\frac{x}{100}+\frac25\cdot\frac{100-x}{100}
=\frac{2x+200}{500}=\frac{x+100}{250},
\]
and that of $B$ is
\[
\Prob_B=\frac{9}{10}\cdot\frac{100-x}{100}+\frac35\cdot\frac{x}{100}
=\frac{900-3x}{1000}.
\]
Hence
\[
\begin{aligned}
\Prob_A>\Prob_B
&\iff \frac{x+100}{250}>\frac{900-3x}{1000}
\iff 4(x+100)>900-3x\\
&\iff 7x>500
\iff x>\frac{500}{7}.
\end{aligned}
\]
So the inequality alone allows every integer $x\ge72$.

However, all row totals must be positive, so $0<x<100$. The success rates
with denominator $5$ require both $x$ and $100-x$ to be multiples of $5$.
In addition, $B$'s easy-task failure count is $(100-x)/10$, so $100-x$
must be a multiple of $10$, and hence so must $x$. Conversely, every
$x\in\{10,20,\ldots,90\}$ gives strictly positive integer entries in both
tables. Together with $x>\frac{500}{7}$, the feasible values are
\[
x\in\{80,\,90\}.
\]

\begin{samepage}
For the smallest such value, $x=80$, the two count tables are
\begin{center}
\begin{tabular}{lccc}
\toprule
$A$ & pass & fail & total\\\midrule
easy & $64$ & $16$ & $80$\\
hard & $8$ & $12$ & $20$\\\midrule
total & $72$ & $28$ & $100$\\\bottomrule
\end{tabular}
\qquad
\begin{tabular}{lccc}
\toprule
$B$ & pass & fail & total\\\midrule
easy & $18$ & $2$ & $20$\\
hard & $48$ & $32$ & $80$\\\midrule
total & $66$ & $34$ & $100$\\\bottomrule
\end{tabular}
\end{center}
Indeed $\Prob_A=\frac{72}{100}=0.72>0.66=\frac{66}{100}=\Prob_B$.
\end{samepage}

\item Both algorithms now receive the same easy-task proportion $t$.
By the law of total probability, conditioning on the task type,
\begin{align*}
\Prob_A &= \frac45 t+\frac25(1-t)=\frac{2+2t}{5},\\
\Prob_B &= \frac{9}{10}t+\frac35(1-t)=\frac{6+3t}{10}.
\end{align*}
Subtracting gives
\[
\Prob_B-\Prob_A
=\frac{6+3t}{10}-\frac{4+4t}{10}
=\frac{2-t}{10}.
\]
Since $0\le t\le1$, we have $2-t\ge1>0$. Thus $B$ necessarily has
the greater overall success probability. The reported difference therefore
gives the equation
\[
\frac{2-t}{10}=\frac{13}{100}
\quad\Longrightarrow\quad
10(2-t)=13
\quad\Longrightarrow\quad
t=\frac7{10}.
\]
Substituting this value into the total-probability formulas,
\begin{align*}
\Prob_A &= \frac45\cdot\frac7{10}+\frac25\cdot\frac3{10}
          =\frac{34}{50}=\frac{17}{25}=0.68,\\
\Prob_B &= \frac{9}{10}\cdot\frac7{10}+\frac35\cdot\frac3{10}
          =\frac{81}{100}=0.81.
\end{align*}
Hence $t=0.7$, algorithm $B$ performs better, and the difference is
$0.81-0.68=0.13$, as required.

\item With easy-task proportions $s$ for $A$ and $t$ for $B$, the law of
total probability gives
\begin{align*}
\Prob_A &= \frac45 s+\frac25(1-s)=\frac{2+2s}{5},\\
\Prob_B &= \frac{9}{10}t+\frac35(1-t)=\frac{6+3t}{10}.
\end{align*}
Requiring $A$'s overall success probability to exceed $B$'s by $1/100$,
\[
\begin{aligned}
\Prob_A-\Prob_B=\frac1{100}
&\iff \frac{4s-3t-2}{10}=\frac1{100}\\
&\iff 4s-3t=\frac{21}{10}
\iff s=\frac{21}{40}+\frac34 t.
\end{aligned}
\]
We next check which values of $t$ satisfy the constraints. Since
$0<t<1$,
\[
s-t=\frac{21}{40}-\frac14 t>\frac{11}{40}>0.
\]
Consequently,
\[
|s-t|\le\frac25
\iff \frac{21}{40}-\frac14 t\le\frac{16}{40}
\iff t\ge\frac12.
\]
Also, $s>0$ holds automatically, while
\[
s<1
\iff \frac{21}{40}+\frac34 t<1
\iff t<\frac{19}{30}.
\]
Combining these conditions, every $t$ in
\[
\frac12\le t<\frac{19}{30},\qquad
s=\frac{21}{40}+\frac34 t,
\]
gives a valid pair. In particular, choose
\[
t=\frac12,\qquad s=\frac9{10}.
\]
Using the law of total probability again, the overall success probabilities
for this pair are
\begin{align*}
\Prob_A &= \frac45\cdot\frac9{10}+\frac25\cdot\frac1{10}
          =\frac{38}{50}=\frac{19}{25}=0.76,\\
\Prob_B &= \frac{9}{10}\cdot\frac12+\frac35\cdot\frac12
          =\frac{15}{20}=\frac34=0.75.
\end{align*}
Both $s=9/10$ and $t=1/2$ lie strictly between $0$ and $1$, so each
algorithm receives both task types. Moreover,
\[
|s-t|=\left|\frac9{10}-\frac12\right|=\frac25,
\qquad
\Prob_A-\Prob_B=\frac{19}{25}-\frac34=\frac1{100}.
\]
Thus all the constraints are satisfied.

This is compatible with $B$'s superiority within each task type because
the overall probabilities are weighted averages with different weights.
Although $B$ is better on both easy and hard tasks, $A$ receives a much
larger proportion of easy tasks ($90\%$ versus $50\%$). In particular,
$A$'s success probability on easy tasks, $4/5$, exceeds $B$'s success
probability on hard tasks, $3/5$. Giving sufficiently more weight to easy
tasks for $A$ can therefore make its overall success probability greater.
\end{enumerate}

\clearpage
\section*{Problem 3}
An automated station is placed in mode $H$ with probability $1/3$ and mode $L$ with probability $2/3$. The mode is chosen once for three checks. Let $A_i$ mean that check $i$ passes. Each check passes with probability $3/4$ conditional on $H$, and with probability $1/4$ conditional on $L$. Let $U=A_1\cap A_2\cap A_3$.
\begin{enumerate}
\item First assume the three pass events are mutually independent conditional on either mode. Compute $\Prob(H\mid U)$. Are $A_1$ and $A_2$ independent without conditioning on the mode? Support your answer with probabilities, and explain the difference between the two independence statements.
\item For the rest of this problem, assume only that every pair of pass events is independent conditional on each mode. For a fixed mode, write $p$ for the common pass probability and $z$ for the probability of three passes. Complete the table below and justify why the outcomes in each grouped row must have equal probabilities. Write $1$ for a pass and $0$ for a failure.
\begin{center}
\begin{tabular}{cc}
\toprule Outcome & Conditional probability of each outcome\\\midrule
$111$ & $z$\\
$110,101,011$ & $\cdots$\\
$100,010,001$ & $\cdots$\\
$000$ & $\cdots$\\\bottomrule
\end{tabular}
\end{center}
Find the full feasible intervals for $h=\Prob(U\mid H)$ and $\ell=\Prob(U\mid L)$. Explain why every pair $(h,\ell)$ in those intervals gives a valid model with the specified marginals and pairwise independence.
\item Find the exact set of possible values of $\Prob(H\mid U)$ under the assumptions of part (b). Prove both that no other value is possible and that every value in your set is attainable. Explain why the single answer in part (a) cannot be inferred from pairwise independence alone.
\end{enumerate}

\subsection*{Solution}

\begin{enumerate}
\item Conditional on either mode, the three pass events are mutually
independent. Therefore,
\[
\Prob(U\mid H)=\left(\frac34\right)^3=\frac{27}{64},
\qquad
\Prob(U\mid L)=\left(\frac14\right)^3=\frac1{64}.
\]
By the law of total probability,
\[
\Prob(U)=\Prob(U\mid H)\Prob(H)+\Prob(U\mid L)\Prob(L)
=\frac{27}{64}\cdot\frac13+\frac1{64}\cdot\frac23
=\frac{29}{192}.
\]
Thus Bayes' formula gives
\[
\Prob(H\mid U)=\frac{\Prob(U\cap H)}{\Prob(U)}
=\frac{\frac{27}{64}\cdot\frac13}{\frac{29}{192}}
=\frac{27}{29}.
\]

For either $i=1$ or $i=2$, the law of total probability gives
\[
\Prob(A_i)=\frac34\cdot\frac13+\frac14\cdot\frac23=\frac5{12}.
\]
Conditional independence also gives
\[
\Prob(A_1\cap A_2)
=\left(\frac34\right)^2\cdot\frac13
 +\left(\frac14\right)^2\cdot\frac23
=\frac{11}{48}.
\]
However,
\[
\Prob(A_1)\Prob(A_2)=\left(\frac5{12}\right)^2
=\frac{25}{144}\ne\frac{33}{144}=\frac{11}{48}.
\]
Hence $A_1$ and $A_2$ are not independent without conditioning on the mode.
Conditional independence describes the checks with the mode fixed.
Unconditional independence concerns the mixture of the two modes.
Here all three checks share the same randomly chosen mode, and observing
one pass provides information about that mode, which changes the
probability of another pass. Thus conditional independence does not imply
unconditional independence.

\item Fix one mode. Throughout this part, probabilities are conditional
on that mode. Write $p$ for each check's pass probability and
$z=\Prob(A_1\cap A_2\cap A_3)$.
Pairwise independence implies
\[
\Prob(A_i\cap A_j)=p^2\qquad(i\ne j).
\]
For example, the event $A_1\cap A_2$ consists of outcomes $111$ and $110$,
so
\[
\Prob(110)=p^2-z.
\]
The same argument using $A_1\cap A_3$ and $A_2\cap A_3$ gives
$\Prob(101)=\Prob(011)=p^2-z$.
Next, using the marginal probability of $A_1$,
\[
p=\Prob(111)+\Prob(110)+\Prob(101)+\Prob(100),
\]
so
\[
\Prob(100)=p-z-2(p^2-z)=p-2p^2+z.
\]
Using the marginals of $A_2$ and $A_3$ in the same way gives the same
probability for $010$ and $001$. Finally, subtracting the other seven
outcome probabilities from $1$,
\[
\Prob(000)=1-z-3(p^2-z)-3(p-2p^2+z)
=1-3p+3p^2-z.
\]
Hence the completed table is
\begin{center}
\begin{tabular}{cc}
\toprule
Outcome & Conditional probability of each outcome\\\midrule
$111$ & $z$\\
$110,101,011$ & $p^2-z$\\
$100,010,001$ & $p-2p^2+z$\\
$000$ & $1-3p+3p^2-z$\\\bottomrule
\end{tabular}
\end{center}
The equal probabilities within each grouped row follow from the equal
marginals and pairwise independence, as derived above.

All eight outcome probabilities must be nonnegative. Thus
\[
z\ge0,\qquad z\le p^2,\qquad
z\ge2p^2-p,\qquad z\le1-3p+3p^2.
\]
Equivalently, the full feasible interval is
\[
\max\{0,2p^2-p\}\le z\le\min\{p^2,1-3p+3p^2\}.
\]
For mode $H$, substitute $p=3/4$:
\[
\max\left\{0,\frac38\right\}\le h
\le\min\left\{\frac9{16},\frac7{16}\right\},
\qquad\text{so}\qquad
\boxed{\frac38\le h\le\frac7{16}}.
\]
For mode $L$, substitute $p=1/4$:
\[
\max\left\{0,-\frac18\right\}\le\ell
\le\min\left\{\frac1{16},\frac7{16}\right\},
\qquad\text{so}\qquad
\boxed{0\le\ell\le\frac1{16}}.
\]

These conditions are also sufficient. For any $z$ in the corresponding
interval, assign the eight outcomes the probabilities in the completed
table. They are nonnegative and sum to
\[
z+3(p^2-z)+3(p-2p^2+z)+(1-3p+3p^2-z)=1.
\]
Each pass marginal is
\[
z+2(p^2-z)+(p-2p^2+z)=p,
\]
and each pair of passes has probability
\[
z+(p^2-z)=p^2,
\]
which proves pairwise independence. Therefore, for every
$(h,\ell)\in[3/8,7/16]\times[0,1/16]$, we can use the table with
$(p,z)=(3/4,h)$ in mode $H$ and $(p,z)=(1/4,\ell)$ in mode $L$.
Choosing $H$ with probability $1/3$ and $L$ with probability $2/3$
then defines a valid model satisfying all the assumptions of part (b).

\item By the law of total probability and the notation from part (b),
\[
\Prob(U)=\Prob(U\mid H)\Prob(H)+\Prob(U\mid L)\Prob(L)
=\frac{h}{3}+\frac{2\ell}{3}.
\]
Since $h\ge3/8>0$, we have $\Prob(U)>0$. Bayes' formula therefore gives
\[
\Prob(H\mid U)
=\frac{\Prob(U\mid H)\Prob(H)}{\Prob(U)}
=\frac{h/3}{h/3+2\ell/3}
=\frac{h}{h+2\ell}
=\frac{1}{1+2\ell/h}.
\]
From part (b), $3/8\le h\le7/16$ and $0\le\ell\le1/16$, so
\[
0\le\frac{\ell}{h}\le\frac{1/16}{3/8}=\frac16.
\]
It follows that
\[
\frac34=\frac{1}{1+2(1/6)}
\le\Prob(H\mid U)\le1.
\]
The lower bound is attained at $(h,\ell)=(3/8,1/16)$, and the upper
bound is attained whenever $\ell=0$.

To prove that every intermediate value is attainable, let
$q\in[3/4,1]$ be arbitrary and choose
\[
h=\frac38,\qquad \ell=\frac{3(1-q)}{16q}.
\]
Because $q\le1$ and $q\ge3/4$, we have
\[
0\le\ell=\frac3{16}\left(\frac1q-1\right)
\le\frac3{16}\left(\frac43-1\right)=\frac1{16}.
\]
Thus this pair is feasible by part (b), and it gives
\[
\Prob(H\mid U)=\frac{3/8}{3/8+3(1-q)/(8q)}=q.
\]
Consequently, the exact set of possible values is
\[
\boxed{\Prob(H\mid U)\in\left[\frac34,1\right]}.
\]
Under the mutual conditional independence assumption of part (a),
$h=(3/4)^3=27/64$ and $\ell=(1/4)^3=1/64$, which give the single value
$27/29$. Pairwise conditional independence fixes the probabilities of
two passes, but does not fix the probability of three passes.
It therefore permits the different feasible pairs $(h,\ell)$ from
part (b), so the single answer $27/29$ cannot be inferred from pairwise
independence alone.
\end{enumerate}

\clearpage
\section*{Problem 4}
In a three-door game, the car is behind doors $1,2,3$ with probabilities
$1/2,1/3,1/6$. You choose door $1$. The host knows the car's location and
opens a goat door from doors $2,3$. If both doors contain goats, host A
opens door $2$ with probability $1/4$, whereas host B does so with
probability $3/4$. Otherwise, the opening is forced.

A fair coin selects one host, independently of the car. The selected host
first performs a calibration round, with the car known to be behind door
$1$ and door $1$ initially selected, and then the main game. Given the host, the two rounds use independent
random choices. Let $C$ and $J$ be the calibration and main openings,
respectively, and let $K$ be the main car door.

System F retains exactly the runs with $C=2$; system M retains exactly
the runs with $C=J$. Write $F=\{C=2\}$ and $M=\{C=J\}$. In a retained
run, the system and $J$ are reported, but the host and $C$ are not
separately reported. All questions below concern the main game.

\begin{enumerate}
\item Prove that the two systems give the same
posterior probability of $K=1$ when $J=2$ is reported, and calculate
this probability. Identify the conditioning events before calculating.

\item For each system, make a two-row table, indexed
by $j=2,3$, giving
\[
\Prob(J=j\mid\text{retained}),\qquad
\Prob(K=1\mid J=j,\text{retained}).
\]
Reuse your result from part (a). Apply the conditional law of total
probability to find \mbox{$\Prob(K=1\mid F)$} and \mbox{$\Prob(K=1\mid M)$}.
Explain why the equality in part (a) does not imply equality of these
two probabilities.

\item Modify only the host selection: now select a
host independently by a fair coin for each of the two rounds. Conditional on the selected hosts, the rounds use independent random choices. Both
recording systems keep their original rules. Does either system change
the joint distribution of $(K,J)$ by retaining runs? Prove your answer
for every pair $(k,j)$ with positive probability. Begin by finding the
probability of retention conditional on $K=k$ and $J=j$, using the law of
total probability over the calibration host. Explain what feature of
the original experiment allowed the two systems to give different
answers in part (b).
\end{enumerate}

\subsection*{Solution}

Let $D\in\{A,B\}$ denote the selected host, with
$\Prob(D=A)=\Prob(D=B)=1/2$. Write $q_A=1/4$ and $q_B=3/4$ for
the probabilities of opening door $2$ when both available doors contain
goats. For each host $d$,
\[
\Prob(C=2\mid D=d)=q_d,\qquad
\Prob(J=2\mid K=1,D=d)=q_d.
\]
If $K=2$, then $J=3$ is forced; if $K=3$, then $J=2$ is forced.
Given the host, the two rounds use independent random choices.

\begin{enumerate}
\item The conditioning events are
\[
F\cap\{J=2\}=\{C=2,J=2\},\qquad
M\cap\{J=2\}=\{C=2,J=2\}.
\]
They are identical, so the two posterior probabilities are equal.
To calculate their common value, condition on the host. Since the host
selection is independent of $K$,
\begin{align*}
\Prob(K=1,C=2,J=2)
&=\sum_{d\in\{A,B\}}\Prob(D=d)\Prob(K=1)q_d^2\\
&=\frac12\cdot\frac12\left[\left(\frac14\right)^2
  +\left(\frac34\right)^2\right]
=\frac5{32}.
\end{align*}
The only other possible car location when $J=2$ is $K=3$, for which the
main opening is forced. Thus
\[
\Prob(K=3,C=2,J=2)
=\frac12\cdot\frac16\left(\frac14+\frac34\right)=\frac1{12}.
\]
Consequently,
\[
\Prob(C=2,J=2)=\frac5{32}+\frac1{12}=\frac{23}{96},
\]
and
\[
\boxed{\Prob(K=1\mid J=2,F)=\Prob(K=1\mid J=2,M)
=\frac{5/32}{23/96}=\frac{15}{23}}.
\]

\item \textbf{System F.} Here retention means $C=2$. For $J=2$, part (a)
gives $\Prob(F,J=2)=23/96$ and the posterior $15/23$.
For $J=3$, the possible car locations are $K=1$ and $K=2$.
Conditioning on the host gives
\begin{align*}
\Prob(K=1,C=2,J=3)
&=\sum_{d\in\{A,B\}}\Prob(D=d)\Prob(K=1)q_d(1-q_d)\\
&=\frac12\cdot\frac12
\left(\frac14\cdot\frac34+\frac34\cdot\frac14\right)
=\frac3{32},\\
\Prob(K=2,C=2,J=3)
&=\frac12\cdot\frac13\left(\frac14+\frac34\right)=\frac16.
\end{align*}
Hence
\[
\Prob(F,J=3)=\frac3{32}+\frac16=\frac{25}{96},\qquad
\Prob(F)=\frac{23}{96}+\frac{25}{96}=\frac12.
\]
Normalizing these joint probabilities gives
\[
\Prob(J=2\mid F)=\frac{23/96}{1/2}=\frac{23}{48},\qquad
\Prob(J=3\mid F)=\frac{25/96}{1/2}=\frac{25}{48},
\]
and
\[
\Prob(K=1\mid J=3,F)=\frac{3/32}{25/96}=\frac9{25}.
\]
The requested table for system F is
\begin{center}
\begin{tabular}{ccc}
\toprule
$j$ & $\Prob(J=j\mid F)$ & $\Prob(K=1\mid J=j,F)$\\\midrule
$2$ & $23/48$ & $15/23$\\
$3$ & $25/48$ & $9/25$\\\bottomrule
\end{tabular}
\end{center}
By the conditional law of total probability,
\[
\Prob(K=1\mid F)
=\sum_{j=2,3}\Prob(K=1\mid J=j,F)\Prob(J=j\mid F)
=\frac{15}{23}\cdot\frac{23}{48}+\frac9{25}\cdot\frac{25}{48}
=\boxed{\frac12}.
\]

\textbf{System M.} Here retention means $C=J$. For $J=2$, part (a)
again gives $\Prob(M,J=2)=23/96$ and the posterior $15/23$.
For $J=3$, retention requires $C=3$. Conditioning on the host,
\begin{align*}
\Prob(K=1,C=3,J=3)
&=\sum_{d\in\{A,B\}}\Prob(D=d)\Prob(K=1)(1-q_d)^2\\
&=\frac12\cdot\frac12\left[\left(\frac34\right)^2
  +\left(\frac14\right)^2\right]
=\frac5{32},\\
\Prob(K=2,C=3,J=3)
&=\frac12\cdot\frac13\left(\frac34+\frac14\right)=\frac16.
\end{align*}
Therefore,
\[
\Prob(M,J=3)=\frac5{32}+\frac16=\frac{31}{96},\qquad
\Prob(M)=\frac{23}{96}+\frac{31}{96}=\frac9{16}.
\]
It follows that
\[
\Prob(J=2\mid M)=\frac{23/96}{9/16}=\frac{23}{54},\qquad
\Prob(J=3\mid M)=\frac{31/96}{9/16}=\frac{31}{54},
\]
and
\[
\Prob(K=1\mid J=3,M)=\frac{5/32}{31/96}=\frac{15}{31}.
\]
The requested table for system M is
\begin{center}
\begin{tabular}{ccc}
\toprule
$j$ & $\Prob(J=j\mid M)$ & $\Prob(K=1\mid J=j,M)$\\\midrule
$2$ & $23/54$ & $15/23$\\
$3$ & $31/54$ & $15/31$\\\bottomrule
\end{tabular}
\end{center}
Applying the conditional law of total probability again,
\[
\Prob(K=1\mid M)
=\sum_{j=2,3}\Prob(K=1\mid J=j,M)\Prob(J=j\mid M)
=\frac{15}{23}\cdot\frac{23}{54}+\frac{15}{31}\cdot\frac{31}{54}
=\boxed{\frac59}.
\]
The equality in part (a) applies only to the row $J=2$. When $J=3$,
system F conditions on $C=2$, while system M conditions on $C=3$, giving
different posterior probabilities. The two systems also assign different
conditional probabilities to $J=2$ and $J=3$. Thus their weighted averages
need not be equal, and here $\Prob(K=1\mid F)=1/2\ne5/9=\Prob(K=1\mid M)$.

\item Let $D_C$ and $D_J$ denote the hosts selected for the calibration
and main rounds, respectively. They are now selected independently, each
by a fair coin, and independently of $K$. Conditional on the selected
hosts, the random choices in the two rounds are independent. Hence the
calibration outcome $C$ is independent of the main pair $(K,J)$.

Fix any $(k,j)$ such that $\Prob(K=k,J=j)>0$. By the law of total
probability over the calibration host,
\begin{align*}
\Prob(F\mid K=k,J=j)
&=\sum_{d\in\{A,B\}}
\Prob(C=2\mid D_C=d,K=k,J=j)\Prob(D_C=d\mid K=k,J=j)\\
&=\frac14\cdot\frac12+\frac34\cdot\frac12
=\frac12.
\end{align*}
For system M, once $J=j$ is given, retention means $C=j$. Thus
\[
\Prob(M\mid K=k,J=j)
=\sum_{d\in\{A,B\}}
\Prob(C=j\mid D_C=d,K=k,J=j)\Prob(D_C=d\mid K=k,J=j).
\]
If $j=2$, this is
\[
\frac14\cdot\frac12+\frac34\cdot\frac12=\frac12;
\]
if $j=3$, it is
\[
\frac34\cdot\frac12+\frac14\cdot\frac12=\frac12.
\]
Therefore, under either system the conditional probability of retention
is $1/2$ for every possible main outcome.

Let $R$ stand for either retention event $F$ or $M$. Summing over all
possible main outcomes gives
\[
\Prob(R)=\sum_{(k,j):\,\Prob(K=k,J=j)>0}
\Prob(R\mid K=k,J=j)\Prob(K=k,J=j)=\frac12.
\]
It follows that, for every pair with positive probability,
\begin{align*}
\Prob(K=k,J=j\mid R)
&=\frac{\Prob(R\mid K=k,J=j)\Prob(K=k,J=j)}{\Prob(R)}\\
&=\frac{(1/2)\Prob(K=k,J=j)}{1/2}
=\Prob(K=k,J=j).
\end{align*}
Thus neither system changes the joint distribution of $(K,J)$ by
retaining runs.

For completeness, averaging over the main host gives
\[
\Prob(J=2\mid K=1)=\Prob(J=3\mid K=1)=\frac12.
\]
The other main openings are forced, so the joint probabilities before
and after retention are
\begin{center}
\begin{tabular}{cccc}
\toprule
$(k,j)$ & $\Prob(K=k,J=j)$ & $\Prob(K=k,J=j\mid F)$
& $\Prob(K=k,J=j\mid M)$\\\midrule
$(1,2)$ & $1/4$ & $1/4$ & $1/4$\\
$(1,3)$ & $1/4$ & $1/4$ & $1/4$\\
$(2,3)$ & $1/3$ & $1/3$ & $1/3$\\
$(3,2)$ & $1/6$ & $1/6$ & $1/6$\\\bottomrule
\end{tabular}
\end{center}
All other pairs have probability zero and remain impossible after
retention.

In the original experiment, the same randomly selected host was used in
both rounds. This shared host made the calibration opening informative
about the host in the main game, linking $C$ to $J$ even though the
rounds used independent random choices conditional on the host.
The two retention rules could therefore favor different main outcomes,
producing the different answers in part (b). With independent host
selections, $C$ is independent of $(K,J)$ and is equally likely to be
$2$ or $3$, so both rules retain every possible main outcome with the
same probability $1/2$.
\end{enumerate}

\clearpage
\section*{Problem 5}
Start with one amoeba at time $0$. At each minute's end, every living
amoeba dies or splits into two. The death probability is $1/4$ in odd
minutes and $p\in(0,1)$ in even minutes. All choices are independent
across individuals and minutes; an empty population stays empty.
Do not use general branching-process theorems.

\begin{enumerate}
\item For integers $n\geq0$, let $x_n,y_n$ be the probabilities of
extinction within $n$ updates, starting with one amoeba under the
alternating rules, with the odd/even rule first, respectively. Let $x,y$ be the corresponding
eventual probabilities. Find the initial values, first-step recursions
and limiting equations. Explain why $x=y$ cannot be assumed.
You may use \mbox{$\Prob(\bigcup_{n\geq0}E_n)=\lim_{n\to\infty}\Prob(E_n)$}
whenever $E_n\subseteq E_{n+1}$ for all $n\geq0$.

\item Set $p=2/3$. Prove $x=y=1$. An averaged model replaces both
death probabilities by their average $\bar p=11/24$, keeping all other rules unchanged;
its eventual extinction probability is given as $11/13$.
Calculate extinction probabilities by minute $2$ in both models.
Does averaging overestimate or underestimate extinction at minute $2$
and eventually?

\item Now $p$ is unknown, and the probability of exactly two amoebae
at the end of minute $2$ is $1/3$. Given exactly four at the end of
minute $3$, find every possible probability that all four descend
from the same one of the two minute-$1$ amoebae, with the corresponding
$p$. Is this conditional probability uniquely determined?
\end{enumerate}

\subsection*{Solution}

\begin{enumerate}
\item At time $0$ there is one living amoeba, so extinction has not yet
occurred. Hence
\[
x_0=y_0=0.
\]
Starting with the odd-minute rule, the first amoeba dies with probability
$1/4$, giving immediate extinction. With probability $3/4$ it splits
into two. Each descendant then starts a process with the even-minute
rule first and has probability $y_n$ of extinction within the remaining
$n$ updates. The choices in the two descendant families are independent,
and both families must become extinct. Thus
\[
\boxed{x_{n+1}=\frac14+\frac34 y_n^2},\qquad n\ge0.
\]
Similarly, starting with the even-minute rule gives
\[
\boxed{y_{n+1}=p+(1-p)x_n^2},\qquad n\ge0.
\]
In particular, $x_1=1/4$ and $y_1=p$.

Substituting one first-step recursion into the other gives the relations
between $n$ and $n+2$:
\begin{align*}
x_{n+2}
&=\frac14+\frac34 y_{n+1}^2
=\boxed{\frac14+\frac34\left[p+(1-p)x_n^2\right]^2},\\
y_{n+2}
&=p+(1-p)x_{n+1}^2
=\boxed{p+(1-p)\left[\frac14+\frac34 y_n^2\right]^2},
\qquad n\ge0.
\end{align*}
After two updates the starting phase is restored, so each two-step
recursion involves only its own sequence.

For either starting phase, the event of extinction within $n$ updates
is contained in the event of extinction within $n+1$ updates. Their union
is the event of eventual extinction. Therefore, by the given limit
property,
\[
x=\lim_{n\to\infty}x_n,\qquad
y=\lim_{n\to\infty}y_n.
\]
Taking limits in the first-step recursions yields
\[
\boxed{x=\frac14+\frac34 y^2},\qquad
\boxed{y=p+(1-p)x^2}.
\]
Taking limits in the two-step recursions also gives
\[
x=\frac14+\frac34\left[p+(1-p)x^2\right]^2,
\qquad
y=p+(1-p)\left[\frac14+\frac34 y^2\right]^2.
\]

We cannot assume $x=y$: the two processes start with different death
probabilities and different orders of the alternating rules. More
explicitly, if $x=y=q$, subtracting the two first-step limiting equations
would give
\[
0=\left(p-\frac14\right)(1-q^2).
\]
Thus when $p\ne1/4$, equality would force $q=1$, since $0\le q\le1$.
That would assert certain extinction, which must be proved rather than
assumed. Equality may hold for particular parameters, but it is not a
consequence of the setup alone.

\item Set $p=2/3$. The limiting equations from part (a) become
\[
x=\frac14+\frac34 y^2,\qquad
y=\frac23+\frac13x^2.
\]
Substitute the second equation into the first:
\[
x=\frac14+\frac34\left(\frac23+\frac13x^2\right)^2
=\frac14+\frac{(2+x^2)^2}{12}.
\]
Multiplying by $12$ and rearranging,
\[
12x=3+(2+x^2)^2
\quad\Longrightarrow\quad
x^4+4x^2-12x+7=0.
\]
The polynomial factors as
\[
(x-1)^2(x^2+2x+7)=0.
\]
Since $x^2+2x+7=(x+1)^2+6>0$, the only real solution is $x=1$.
Substituting back gives $y=2/3+(1/3)\cdot1^2=1$. Hence
\[
\boxed{x=y=1}.
\]
This conclusion follows directly from the limiting equations, without
using a general branching-process theorem.

We next calculate extinction by minute $2$, starting with one amoeba
at time $0$ and applying the odd-minute rule first. Extinction occurs
either when the original amoeba dies at minute $1$, or when it splits
at minute $1$ and both descendants die at minute $2$. These cases are
disjoint. By the law of total probability and independence, the
alternating model gives
\[
\Prob(\text{extinction by minute }2)
=\frac14+\frac34\left(\frac23\right)^2
=\frac14+\frac13
=\boxed{\frac7{12}}.
\]
In the averaged model, the death probability at each minute is
$\bar p=11/24$. Applying the same two-case calculation,
\[
\begin{aligned}
\Prob_{\mathrm{avg}}(\text{extinction by minute }2)
&=\bar p+(1-\bar p)\bar p^2\\
&=\frac{11}{24}+\frac{13}{24}\left(\frac{11}{24}\right)^2
=\boxed{\frac{7909}{13824}}.
\end{aligned}
\]
Since
\[
\frac7{12}-\frac{7909}{13824}
=\frac{8064-7909}{13824}
=\frac{155}{13824}>0,
\]
averaging \emph{underestimates} extinction by minute $2$.
For eventual extinction, the averaged probability is given as $11/13$,
whereas the alternating model has probability $x=1$. Therefore,
\[
\frac{11}{13}<1,
\]
so averaging also \emph{underestimates} eventual extinction.

\item Let $N_m$ denote the population at the end of minute $m$.
To have $N_2=2$, the initial amoeba must split at minute $1$, and exactly
one of the two minute-$1$ amoebae must split at minute $2$ while the
other dies. By independence,
\[
\Prob(N_2=2)=\frac34\cdot2p(1-p)=\frac13.
\]
Thus
\[
p(1-p)=\frac29
\quad\Longrightarrow\quad
9p^2-9p+2=(3p-1)(3p-2)=0.
\]
Both roots lie in $(0,1)$, so the possible parameters are
\[
\boxed{p=\frac13\quad\text{or}\quad p=\frac23}.
\]

Write $S=\{N_1=2\}$. The event $N_3=4$ implies $S$, since a population
that dies at minute $1$ cannot return. Conditional on $S$, label the
two minute-$1$ amoebae $1$ and $2$, and let $V_i$ be the number of
living minute-$3$ descendants of amoeba $i$. Their descendant families
evolve independently and have the same distribution.

Each $V_i$ can be $0$, $2$, or $4$. For $V_i=0$, its ancestor either
dies at minute $2$, or splits at minute $2$ and both children die at
minute $3$. For $V_i=2$, it splits at minute $2$ and exactly one child
splits at minute $3$. For $V_i=4$, it splits at minute $2$ and both
children split at minute $3$. Therefore,
\begin{align*}
u_0:=\Prob(V_i=0\mid S)
&=p+(1-p)\left(\frac14\right)^2=\frac{1+15p}{16},\\
u_2:=\Prob(V_i=2\mid S)
&=(1-p)\cdot2\cdot\frac14\cdot\frac34=\frac{3(1-p)}8,\\
u_4:=\Prob(V_i=4\mid S)
&=(1-p)\left(\frac34\right)^2=\frac{9(1-p)}{16}.
\end{align*}
The event $N_3=4$ consists of the three disjoint possibilities
\[
(V_1,V_2)=(0,4),\ (2,2),\ (4,0).
\]
Let $G$ mean that all four living minute-$3$ amoebae descend from the
same minute-$1$ amoeba. Among these possibilities, $G$ holds exactly
for $(0,4)$ and $(4,0)$. Hence
\[
\Prob(G,N_3=4\mid S)=2u_0u_4,\qquad
\Prob(N_3=4\mid S)=2u_0u_4+u_2^2.
\]
Since both events imply $S$, the factor $\Prob(S)=3/4$ cancels in the
conditional probability. Thus
\begin{align*}
\Prob(G\mid N_3=4)
&=\frac{\Prob(S)\,2u_0u_4}
        {\Prob(S)(2u_0u_4+u_2^2)}\\
&=\frac{18(1+15p)(1-p)}
        {18(1+15p)(1-p)+36(1-p)^2}\\
&=\frac{1+15p}{3+13p},
\end{align*}
where the second line follows by multiplying numerator and denominator
by $256$, and the last line uses $p<1$.

For $p=1/3$, this gives
\[
\Prob(G\mid N_3=4)=\frac{1+15/3}{3+13/3}
=\boxed{\frac9{11}}.
\]
For $p=2/3$, it gives
\[
\Prob(G\mid N_3=4)=\frac{1+15(2/3)}{3+13(2/3)}
=\boxed{\frac{33}{35}}.
\]
These are all the possible values, because the given population
probability allows only the two parameters found above, and both yield
$\Prob(N_3=4)>0$. Since $9/11\ne33/35$, the conditional probability is
\emph{not uniquely determined}.
\end{enumerate}

\end{document}
